\chapter{Conclusiones}
\label{chap:conclusiones}

%----------------------------------------------------------------------------------------

\section{Conclusiones Generales}
\label{sec:conclusiones_generales}

La solución implementada consolidó una arquitectura de nivel empresarial para firma electrónica sobre infraestructura institucional de la FACET. El despliegue integró Documenso como núcleo de aplicación, PostgreSQL para persistencia transaccional, Browserless para renderizado del certificado de firma, MinIO como almacenamiento de objetos y el \textit{proxy} Socat para enlazar la red Docker con el servidor físico de almacenamiento, bajo la orquestación de Coolify y Traefik en un contenedor LXC sobre Proxmox VE. Las verificaciones del Capítulo~\ref{chap:ensayos_resultados} ---desde \textit{healthchecks} de servicios hasta ensayos de extremo a extremo y restauración de respaldos--- validaron el cumplimiento de los objetivos funcionales, de seguridad, de rendimiento y de continuidad operativa definidos en la Sección~\ref{sec:objetivos_alcance}, en concordancia con la matriz de trazabilidad de la Subsección~\ref{subsec:trazabilidad_requisitos}.

La adopción de una alternativa de código abierto auditada y desplegada en entorno propio entregó un retorno de inversión concreto: eliminó el costo recurrente de licencias comerciales, preservó la soberanía total de los datos y evitó la dependencia contractual de proveedores externos de \textit{software}. Frente al escenario de suscripción anual proyectado en el Capítulo~\ref{chap:introduccion}, la arquitectura implementada resultó financieramente sustentable para una institución pública, al trasladar la inversión hacia capacidades internas reutilizables de integración, seguridad y operación.

Este resultado constituyó, además, uno de los primeros servicios productivos desplegados íntegramente sobre la infraestructura del DEEC. La operación sostenida confirmó que los recursos disponibles de virtualización, almacenamiento y red ---Proxmox VE, TrueNAS/MinIO y Coolify--- soportan cargas reales con criterios de disponibilidad, seguridad y trazabilidad equivalentes a un entorno operativo formal. En términos institucionales, la implementación establece una base operativa concreta sobre la que la FACET puede desplegar futuros servicios sin partir de cero.

La integración de este ecosistema puso en juego conocimientos transversales de la carrera: la orquestación de la topología requirió aplicar conceptos de sistemas operativos y protocolos de comunicación TCP/IP para segmentar el tráfico mediante túneles y \textit{proxies} inversos; el diseño de la retención de datos se sustentó en principios de bases de datos y continuidad operativa; el diagnóstico de las limitaciones del procesador demostró que la arquitectura de computadoras es determinante en decisiones de compatibilidad; y la evaluación de alternativas y el análisis de impacto económico se apoyaron en criterios de gestión de proyectos y planificación tecnológica.

El cumplimiento del requisito de usabilidad se verificó mediante una arquitectura web \textit{clientless}: los usuarios operan el circuito de firma desde navegadores estándar de escritorio o móviles, sin instalar \textit{plugins} ni \textit{software} de terceros. Esa propiedad eliminó la fricción de adopción y redujo los costos de soporte en puestos finales, dos factores que refuerzan directamente la viabilidad operativa del servicio a escala institucional.

Lo que este proyecto enseñó con mayor claridad es que los registros, las métricas y los respaldos no son un complemento del sistema, sino su columna vertebral. Sin los registros de contenedores no habría sido posible diagnosticar el fallo de sellado de Browserless ni identificar los flujos de acceso a MinIO; sin las métricas de la Sección~\ref{sec:metricas_rendimiento}, los picos de RAM de Browserless bajo carga habrían sido invisibles hasta convertirse en un problema real. Y lo que quedó más claro de todo es que ninguna intervención sobre el entorno productivo es aceptable sin un respaldo verificado previo: la prueba de restauración documentada en la Sección~\ref{sec:validacion_respaldo} no fue un ejercicio opcional, fue la única manera de saber con certeza que los volcados almacenados en MinIO realmente servían para algo.

En cuanto a la continuidad operativa, la estrategia de respaldo automatizado hacia el almacenamiento externo ---junto con los procedimientos de restauración del Apéndice~\ref{AppendixC} y la guía de inicio rápido del Apéndice~\ref{AppendixD}--- garantizan que el sistema pueda ser recuperado y operado por el personal de la facultad sin depender del autor original. Esa transferencia de conocimiento es, en última instancia, el criterio más exigente de cualquier despliegue productivo: que el servicio sobreviva a quien lo construyó.

%----------------------------------------------------------------------------------------

\section{Trabajo Futuro}
\label{sec:trabajo_futuro}

Si bien el sistema cumple con los objetivos planteados y se encuentra operativo en producción, la arquitectura actual admite tres ejes de evolución prioritarios: resiliencia, escalabilidad horizontal y extensión institucional.

\begin{itemize}
    \item Respaldo \textit{off-site} y replicación. Tal como se advirtió en la Sección~3.5.2 del diseño, el servidor TrueNAS actual constituye un punto único de fallo (\textit{SPOF}) para el repositorio documental. La siguiente etapa debe implementar \textit{Cross-Site Replication} en MinIO hacia un nodo geográficamente separado o hacia un proveedor de nube pública, para asegurar capacidad efectiva de \textit{Disaster Recovery} ante incidentes de sitio completo.
    \item Alta Disponibilidad del nodo de cómputo. Si bien el repositorio documental queda salvaguardado mediante la replicación descrita en el punto anterior, el servidor físico actual constituye un \textit{SPOF} a nivel de cómputo: una interrupción del \textit{hardware} del nodo provoca la caída inmediata del servicio. La arquitectura debe evolucionar hacia un clúster de Alta Disponibilidad gestionado mediante \textit{Docker Swarm} o \textit{Kubernetes/K3s}, de modo que la carga se redistribuya ante la falla de un nodo y el servicio mantenga continuidad operativa.
    \item Migración de \textit{hardware} y actualización de versión. El servidor actual utiliza procesadores sin soporte de instrucciones \textit{AVX}, lo que restringe el despliegue a Documenso v2.1.0. La migración a \textit{hardware} compatible habilitará la adopción de ramas superiores con nuevas funcionalidades y parches de seguridad vigentes; el procedimiento operativo de actualización se encuentra documentado en el Apéndice~\ref{AppendixC:ActualizacionVersion}.
    \item Extensión del servicio a otras dependencias de la universidad. El modelo de organizaciones y equipos de Documenso permite incorporar nuevas facultades o áreas de la UNT sobre la misma plataforma, preservando aislamiento lógico por unidad organizativa y evitando duplicación innecesaria de infraestructura.
    \item Obtención de certificado emitido por una autoridad certificante de confianza. La infraestructura desplegada soporta de forma nativa certificados X.509 emitidos por autoridades certificantes licenciadas. La incorporación de un certificado anclado en una cadena de confianza pública permitirá evolucionar de firma electrónica a firma digital conforme a la Ley~25.506, sin rediseñar la arquitectura base.
\end{itemize}
